\section{拉格朗日插值}

插值是一种通过已知的、离散的数据点推算一定范围内的新数据点的方法．插值法常用于函数拟合中．

对已知的 $n$ 个点 $(x_0,y_0),(x_1,y_1),\dots,(x_{n-1},y_{n-1})$，求次数不超过 $n-1$ 的多项式 $f(x)$ 满足 $f(x_i)=y_i$（$i=0,\dots,n-1$）。

Lagrange 插值的形式为：

$$f(x)=\sum_{i=1}^ny_i\cdot\prod_{j\neq i}\dfrac{x-x_j}{x_i-x_j}$$

朴素实现的时间复杂度为 $\mathcal{O}(n^2)$，可以优化到 $\mathcal{O}(n\log^2n)$

\subsubsection{普通版本}

\begin{lstlisting}
template<typename T>
std::vector<T> Lagrange_Interpolation(std::vector<T>& x,std::vector<T>& y){
    int n = x.size();
    std::vector<T> f(n),g(n+1),px(n,1);
    g[0] = 1;
    for (int i = 0;i < n;i ++){
        for (int j = i;j >= 0;j --){
            g[j+1] += g[j];
            g[j] *= -x[i];
        }
    }
    for (int i = 0;i < n;i ++){
        for (int j = 0;j < n;j ++){
            if (i == j) continue;
            px[i] *= x[i] - x[j];
        }
    }
    for (int i = 0;i < n;i ++){
        T c = y[i] / px[i],h = g[n];
        for (int j = n-1;j >= 0;j --){
            f[j] += h * c;
            h = g[j] + h * x[i];
        }
    }
    return f;
}
\end{lstlisting}

\subsubsection{$x$ 是连续整数版本}

优化上述式子，最终其插值式子为：

\begin{equation}\begin{split}f(x)&=\sum_{i=1}^{n}y_i\prod_{j\neq i}\frac{x-j}{i-j}\\ &=\sum_{i=1}^{n}y_i(-1)^{n-i}\frac{pre_{i-1}\cdot suf_{i+1}}{(i-1)!(n-i)!}\end{split}\end{equation}

\begin{lstlisting}
template<typename T>
T Lagrange_Interpolation(std::vector<T> y,T x){ // For x=[1..n]
    int n = y.size();
    std::vector<T> fac(n+1,1),pre(n+2),suf(n+2);
    for (int i = 1;i <= n;i ++) fac[i] = fac[i-1] * i;
    pre[0] = suf[n+1] = 1;
    for (int i = 1;i <= n;i ++) pre[i] = pre[i-1] * (x - i);
    for (int i = n;i >= 1;i --) suf[i] = suf[i+1] * (x - i);
    T res = 0;
    for (int i = 1;i <= n;i ++){
        res += y[i-1] * pre[i-1] * suf[i+1] / (fac[i-1] * fac[n-i]) * ((n-i) % 2 == 0 ? 1 : -1);
    }
    return res;
}
\end{lstlisting}